\chapter{Discussion Approfondie et Analyse Critique}
\label{chap:discussion}

Ce chapitre propose une analyse critique et approfondie des comportements observés lors de nos expérimentations. Nous y décortiquons les mécanismes d'autorégulation de l'essaim, l'interprétation des tests statistiques, le bilan de complexité calculatoire et les enseignements au regard du théorème du \textit{No Free Lunch}.

\section{Autorégulation de l'Essaim par la Dynamique du Réplicateur}

\subsection{Un Régulateur Cybernétique à Rétroaction Empirique}
La principale valeur ajoutée de la Théorie des Jeux Évolutive réside dans sa capacité à remplacer un contrôle pré-programmé par un \textbf{système cybernétique bouclé}. Dans le GWO standard, le paramètre $a$ décroît de façon aveugle :
\begin{equation}
    a(t) = 2 - 2 \cdot \frac{t}{T_{\max}}
\end{equation}
Ce modèle suppose implicitement que le paysage d'optimisation est uniformément lisse et que la meute aura localisé le bassin d'attraction global exactement à mi-parcours ($t = T_{\max}/2$). Lorsque cette hypothèse échoue (cas des fonctions multimodales trompeuses ou des espaces combinatoires discrets), l'algorithme force une exploitation locale précipitée sur des solutions sous-optimales.

À l'opposé, la dynamique du réplicateur :
\begin{equation}
    p_k(t+1) = p_k(t) + \eta \cdot p_k(t) \cdot (F_k(t) - \bar{F}(t))
\end{equation}
mesure en temps réel le succès effectif des opérateurs. Tant que l'essaim stagne, les opérateurs d'exploitation (\texttt{ALPHA\_EXPLOIT} ou \texttt{LDWPSO}) génèrent un gain nul ($F_k = 0$). Le terme $(F_k - \bar{F})$ devient alors négatif, entraînant automatiquement la régression des comportements d'exploitation au profit des opérateurs de perturbation stochastique (\texttt{LEVY\_EXPLORE}, \texttt{SOCIAL\_PEER}, \texttt{LIPS}).

\subsection{Le Rôle Crucial des Sauts de Lévy}
Sur la fonction de Rastrigin, la présence de multiples crêtes et vallées périodiques bloque irrémédiablement les algorithmes à gradient ou à guidage strictement barycentrique. Le vol de Lévy, modélisé par l'algorithme de Mantegna \cite{mantegna1994fast}, introduit une marche aléatoire dont la longueur des pas obéit à une distribution en loi de puissance :
\begin{equation}
    P(s) \sim s^{-(1 + \beta_{\text{levy}})}, \qquad 0 < \beta_{\text{levy}} \le 2
\end{equation}
Cette propriété permet de concilier une exploration minutieuse locale (petits pas fréquents) et des sauts d'échelle macroscopique occasionnels mais suffisants pour franchir les barrières de potentiel séparant deux bassins d'attraction.

% -----------------------------------------------------------------------------
\section{Interprétation Statistique et Robustesse des Résultats}

Dans le protocole continu sur Rastrigin, le test de Wilcoxon apparié entre GWO classique et GT-GWO indique une $p$-value de $0.317$. Cette valeur, supérieure au seuil usuel de $0.05$, mérite une analyse rigoureuse :
\begin{itemize}[label=\textbullet]
    \item Sur 30 runs, dans 29 cas sur 30, le GWO classique et GT-GWO convergent tous deux vers le zéro machine parfait ($f = 0.0$). Pour ces 29 runs appariés, la différence $d_i = f_{\text{GWO}}^{(i)} - f_{\text{GT-GWO}}^{(i)}$ est exactement égale à 0.
    \item Or, le test de rang signé de Wilcoxon élimine par construction les paires d'exécutions identiques ($d_i = 0$) du calcul des rangs. Le test ne s'exécute donc mathématiquement que sur le seul run discordant (run n$^\circ$10), où GWO stagne à $0.1479$ alors que GT-GWO converge à $0.0$.
    \item Sur un échantillon résiduel d'une seule observation non nulle, la statistique de rang ne peut rejeter formellement l'hypothèse nulle à un niveau de confiance de $5\%$.
\end{itemize}
Cependant, du point de vue de l'ingénierie et de la sûreté de fonctionnement, le critère déterminant est la \textbf{fiabilité opérationnelle} : GT-GWO garantit un taux d'échec nul ($0\%$) contre $3.3\%$ pour le GWO classique, éliminant totalement les risques de divergence critique.

% -----------------------------------------------------------------------------
\section{Analyse de Complexité et Surcoût Calculatoire}

\subsection{Complexité Temporelle et Spatiale}
L'intégration de la Théorie des Jeux Évolutive engendre un surcoût algorithmique minime :
\begin{itemize}[label=\textbullet]
    \item \textbf{Complexité spatiale :} Le stockage se limite aux probabilités $p_k$, aux compteurs $N_k$ et aux gains cumulés des $K = 4$ stratégies. L'empreinte mémoire additionnelle est en $\mathcal{O}(K) = \mathcal{O}(1)$, ce qui est négligeable devant la mémoire requise pour la population $\mathcal{O}(N \cdot d)$.
    \item \textbf{Complexité temporelle :} À chaque itération, l'évaluation du réplicateur nécessite un balayage des $K$ stratégies, soit $\mathcal{O}(K)$ opérations élémentaires. Cette charge est infinitésimale par rapport aux $N$ évaluations de la fonction objectif à $d$ dimensions.
\end{itemize}

\subsection{Le Paradoxe de la Vitesse en Espace Discret}
Dans l'expérimentation du Sac à Dos 0/1, l'algorithme hybride BGT-GWO s'est révélé \textbf{1.67 fois plus rapide} que le BGWO classique ($1.11$ s contre $1.86$ s). Ce résultat, a priori contre-intuitif au vu de la complexité théorique, s'explique par la dynamique combinatoire :
\begin{enumerate}[label=\textbf{\arabic*.}]
    \item Le BGWO standard génère fréquemment des solutions binaires infaisables en raison de la dérive stochastique continue non contrôlée ;
    \item Chaque violation de la capacité contraint l'algorithme à invoquer l'opérateur de réparation glouton, qui requiert un tri partiel de complexité $\mathcal{O}(N_{\text{items}} \log N_{\text{items}})$ ;
    \item BGT-GWO, grâce à l'apprentissage entre pairs et à l'arbitrage EGT, stabilise plus rapidement ses vecteurs continus dans des hyperplans compatibles avec la contrainte de capacité. Le nombre d'objets à éliminer lors des réparations est drastiquement réduit, accélérant directement le temps global d'exécution CPU.
\end{enumerate}

% -----------------------------------------------------------------------------
\section{Limites Observées et Théorème du "No Free Lunch"}

L'expérimentation comparant PSO et SDPSO sur la suite de six benchmarks continus révèle que le PSO classique obtient un fitness inférieur à court terme sur les fonctions unimodales (Sphere, Schwefel 2.22). Cette observation illustre parfaitement le célèbre \textbf{Théorème du No Free Lunch} (NFL) énoncé par Wolpert \& Macready en 1997 \cite{wolpert1997no} : aucun algorithme d'optimisation ne peut surpasser tous les autres sur l'ensemble de tous les problèmes possibles.

Dans notre cas, SDPSO réserve structurellement 30\% de sa population à l'exploration via CLPSO et commute fréquemment vers LIPS et UPSO. Sur une fonction unimodale pure comme Sphere (où il n'existe aucun minimum local trompeur), toute exploration est superflue : la stratégie optimale consiste à exploiter de façon purement opportuniste vers le centre. SDPSO paie le prix de sa prudence diversificatrice sous un budget d'évaluations restreint (3000 FES). Néanmoins, dès que le paysage devient fortement multimodal ou décorrélé, cette réserve d'exploration devient le seul garant contre la stagnation prématurée.
